\section{群论统一：从扫描半群到稳定类型半环}\label{sec:group-unification}

\subsection{迭代作为“数”的生成}

单生成半群 $\langle g\rangle$ 的元素 $g^t$ 可被解释为“时间步”或“迭代次数”的抽象对象。通过前述同构，可将这些迭代对象在稳定地址空间中赋予自然数语义：
\begin{itemize}
  \item “数”不作为先验集合元素出现，而作为扫描迭代的等价类与其稳定编码出现；
  \item “加法”对应迭代的复合在编码空间的读出；
  \item “乘法”对应迭代复合的二层展开在编码空间的读出。
\end{itemize}

\subsection{黄金比例的代数后果：递归尺度与自相似}

Fibonacci 递归由 $\varphi$ 的特征方程 $x^2=x+1$ 统摄。稳定类型基数 $\abs{X_m}=F_{m+2}$ 的渐近增长率为 $\varphi^m$，折叠多重度的增长率为 $(2/\varphi)^m$。这给出一个严格的“尺度流”模型：分辨率提高（$m\to m+1$）时稳定类型空间自相似增长，但微态空间增长更快，从而必然需要折叠；折叠的统计结构由 $\varphi$ 决定。

\paragraph{分辨率算术的统一窗（域相）}
由定理 \ref{thm:finite-resolution-mod}，分辨率 $m$ 的稳定类型算术满足
\((X_m,\boxplus_m,\boxtimes_m)\cong \ZZ/(F_{m+2}\ZZ)\)。
因此“统一/破缺”首先是模数的环论性质：当 $F_{m+2}$ 为素数时进入域相（推论 \ref{cor:field-phase-fib-prime}），不会产生非平凡 CRT 直积分裂；
当 $F_{m+2}$ 含多个不同素因子时，CRT 给出规范的多因子分裂模板（推论 \ref{cor:crt-factorization}）。
表 \ref{tab:fib_prime_resolution_windows} 列出一组 Fibonacci-prime 分辨率窗口（域相“壳层”）及其尺度 $\varepsilon_m=\varphi^{-m}$。
\input{sections/generated/tab_fib_prime_resolution_windows}

\paragraph{最小多轴统一窗（2/3/5 轴的 CRT 并存）}
若把“可观测的离散同余轴”最小化为 \(\bmod 2,\bmod 3,\bmod 5\) 三条基本轴，
则在模结构同构下其并存读出等价于
\(
2,3,5\mid F_{m+2}
\)
（从而 CRT 投影到 \(\ZZ/2\ZZ,\ZZ/3\ZZ,\ZZ/5\ZZ\) 同时存在）。
命题 \ref{prop:crt-235-min-depth} 给出该条件的最小解为 \(m_{\min}=58\)（即 \(m_{\min}+2=60\)），
并由注 \ref{rem:crt-235-m58-factorization} 可见该最小窗的模数 \(F_{60}\) 同时携带一组额外素因子分量。
因此“算术统一深度”可被组织为一个纯环论的选择律：先由三轴并存强制出最小候选窗，
再在该窗及其同余子塔上用扭曲谱/碰撞谱等可审计指纹判定进一步的“对齐/破缺”结构。

\begin{remark}[三轴最小窗上的折叠碰撞耦合量级]\label{rem:gut-m58-collision-gap-scale}
附录 \ref{app:fold-multiplicity} 已将折叠输出的均匀性缺口记为 \(g_{\mathrm{col}}^2(m)=\mathsf{Col}_m-1/F_{m+2}\)。
在最小三轴窗 \(m=58\) 上，由三态碰撞核 \(A_2\) 的精确递推（命题 \ref{prop:pom-s2-recurrence}）可直接计算得到
\input{sections/generated/eq_gut_crt_235_m58_collision_gap}
\par
另一方面，推论 \ref{cor:fold-col-gap-exp-upper} 给出对所有 \(m\ge 2\) 的显式指数上界
\(
g_{\mathrm{col}}^2(m)\le \mathsf{Col}_m\le C_2\,(r_2/4)^m
\)
（其中 \(r_2=\rho(A_2)\)、\(C_2=2(\mathbf{1}^{\mathsf T}v)/r_2^2\)），
从而“深窗 \(\Rightarrow\) 弱耦合”可被完全归约为可计算的 Perron 常数与有限维核读出。
\end{remark}

\subsection{稳定状态作为“黄金本征态”}

考虑黄金均值约束对应的邻接矩阵
\[
A=\begin{pmatrix}
1 & 1\\
1 & 0
\end{pmatrix},
\]
其谱半径为 $\rho(A)=\varphi$。由 Perron--Frobenius 理论得到唯一正特征向量（归一化后给出平稳分布/最大熵马尔可夫测度）。在符号动力系统语境中，该平稳测度是对语法约束最稳定的统计状态，可作为“稳定本征态”的纯数学对应物：
其显式形式（Parry 最大熵测度）见注 \ref{prop:parry-measure-golden}。

\begin{remark}[黄金均值 SFT 的 Parry 基线]\label{prop:parry-measure-golden}
黄金均值 shift 的 Parry 最大熵测度为两状态马尔可夫链，其转移概率为
\[
P_{00}=\varphi^{-1},\quad P_{01}=\varphi^{-2},\quad P_{10}=1.
\]
其平稳分布为
\[
\pi(0)=\frac{\varphi^2}{\varphi^2+1},\qquad
\pi(1)=\frac{1}{\varphi^2+1}.
\]
参见 \cite{LindMarcus1995}。
\end{remark}

由此可把该测度在长度 $m$ 的稳定词空间 $X_m$ 上具体写为：对 $u=u_1\cdots u_m\in X_m$，
\[
\pi_m(u)=\pi(u_1)\prod_{j=1}^{m-1} P_{u_j u_{j+1}}.
\]
注 \ref{prop:parry-measure-golden} 因而为第 \ref{sec:quant-fit} 节的“拟合准则”提供一个封闭基线；其推导与一般有限型子移位完全一致，见 \cite{LindMarcus1995}。

\input{sections/08_group_unification_parry_endpoint_bridge.tex}

\subsection{边界词塔：SM 的分辨率分解与 Zeckendorf--Fibonacci 的 GUT 签名}\label{subsec:bdry-tower-zeck-gut}

\paragraph{端点 \texorpdfstring{$(1,1)$}{(1,1)} 纤维作为“边界词”层}
在黄金均值语法（禁止相邻 \(11\)）下，长度 \(m\) 的稳定类型集合为 \(X_m\)，并且端点纤维
\[
X_m^{\mathrm{bdry}}
:=\{u\in X_m:\ u_1=1,\ u_m=1\}
\]
的基数由 Fibonacci 数强制给出。具体地，推论 \ref{cor:parry-golden-three-levels} 给出
\[
|X_m^{\mathrm{bdry}}|
=\#\{u\in X_m:u_1=1,u_m=1\}
=F_{m-2}.
\]
因此每个分辨率 \(m\) 都携带一个离散的“边界模态数”\(d(m):=|X_m^{\mathrm{bdry}}|\)。

\paragraph{新结论 A（边界词三层塔）：SM 的 \(12\) 个规范玻色子}
在 \(m=6\) 的锚定接口上，\(X_6\) 的 \(18\oplus 3\) 分解提供了三条边界连接类占位符；
若把 gauge 自由度从“单层占位”提升为“跨分辨率塔状模态”，则 SM 的玻色子数
\(\,1+3+8=12\,\)
在本体系中出现为一个严格的 Fibonacci 边界计数分解：
\input{sections/generated/eq_boundary_tower_sm12}
这一步并非重复“\(SU(3)\times SU(2)\times U(1)\) 三因子”结论，而是把接口升级到\emph{玻色子计数层级}：
\(U(1)\)、\(SU(2)\)、\(SU(3)\) 的规范玻色子分别位于 \(m=4,6,8\) 三个边界层，
并以边界词本身作为离散标签候选集。

\begin{remark}[色对称的闭层包络来源：三拷贝 Hilbert 空间的基混合冗余]\label{rem:color-as-three-copy-envelope}
上段给出的 \(m=8\) 边界层为“色”玻色子计数提供了离散锚点。
若在闭层机制上进一步满足：内部态空间可分解为三份不相交拷贝
\[
\mathcal{H}=\mathcal{H}_q\oplus\mathcal{H}_q\oplus\mathcal{H}_q,
\]
且可观测统计对这三份拷贝的基变换不敏感，
则“拷贝混合”的冗余群自然为 \(U(3)\)；在去掉总体相位后得到连续包络的 \(SU(3)\)。
在本文的编译口径中，类似的“拷贝并置”已以类型提升的形式出现（例如 moment-kernel 的 \(\mathrm{LIFT}_{G^{\times q}}\) 并置提升），
从而把连续包络对称的来源落回到一个可追踪的离散拷贝结构上。
\end{remark}

\paragraph{新结论 C（Fold6 的 \texorpdfstring{$18\oplus 3$}{18+3} 锚定强迫一个 \texorpdfstring{$B_3/C_3$}{B3/C3} 统一种子）}
在 \(m=6\) 的锚定接口上，除去 \(X_6=X_6^{\mathrm{cyc}}\sqcup X_6^{\mathrm{bdry}}\) 的 \(18\oplus 3\) 分解外，
\(X_6^{\mathrm{cyc}}\) 内部还出现一个严格的 \(6+12\) 二层指纹（按 Hamming weight 分层）：
\input{sections/generated/eq_fold6_b3c3_seed}
若把三条端点纤维 \(X_6^{\mathrm{bdry}}\) 理解为三条对易连接方向（Cartan-like），把 \(18\) 个闭词 \(X_6^{\mathrm{cyc}}\) 理解为根方向（charged/root-like），
则 \((3+18)=21\) 的计数与 rank-\(3\) 简单李代数的 Cartan/根分解同型，从而把候选统一李型锁死到 \(B_3\) 或 \(C_3\)（两者互为长短根交换对偶类型）。
更强地，\(6+12\) 的 cyclic 指纹与 \(B_3/C_3\) 根系的“两根长”结构一致：一类长度 \(6\)、另一类长度 \(12\)。
\par
对应的离散对称核为 Weyl 群
\[
W(B_3)\cong W(C_3)\cong (\ZZ_2)^3\rtimes S_3,\qquad |W|=48,
\]
其中 \(S_3\) 可读作三通道/三轴的置换，而 \((\ZZ_2)^3\) 对应每个对易方向上的符号反射（电荷镜像）。
\par
为避免“手工配对”，脚本 \texttt{scripts/exp\_fold6\_b3c3\_seed.py} 输出两种规范化（长短根互换）的根字典表（表 \ref{tab:fold6_b3c3_root_dictionary_b3} 与表 \ref{tab:fold6_b3c3_root_dictionary_c3}；见附录），并写出 JSON 证书 \texttt{artifacts/export/fold6\_b3c3\_seed.json}。

\begin{remark}[三通道势—响应辛相空间与 \texorpdfstring{$C_3=\mathrm{Sp}(6)$}{C3=Sp(6)} 连续包络]\label{rem:sp6-minimal-continuous-envelope}
在附录 \ref{app:vector-potential} 的向量势口径下，给定三维势参数 \(\theta=(\theta_1,\theta_2,\theta_3)\) 的压力
\(
P(\theta)=\log\rho(B(\theta))
\)
可完全由有限矩阵的谱半径读出。
对每一通道 \(i\in\{1,2,3\}\) 定义响应（密度/流）
\[
J_i(\theta):=\partial_{\theta_i}P(\theta),
\]
则 \((\theta;J)\in\RR^6\) 给出一个自然的“势—响应相空间”，并携带标准辛形式
\[
\omega:=\sum_{i=1}^{3}\dd\theta_i\wedge \dd J_i.
\]
允许三通道混合但保持共轭结构（即保持 \(\omega\)）的最小连续线性群为
\[
\mathrm{Aut}(\omega)=\mathrm{Sp}(6),
\]
它在本文的 \(B_3/C_3\) 离散统一种子中对应 \(C_3\)（长短根对偶类型的“辛侧”）。
更重要的是，该连续包络与 \(m=6\) 的离散锚点在生成元计数上精确对齐：
\[
\dim\mathfrak{sp}(6)=3(2\cdot 3+1)=21=|X_6|.
\]
因此，把 \(m=6\) 的 \(21\) 维计数接口解释为 rank-\(3\) 统一代数的最小连续包络时，\(\mathrm{Sp}(6)\) 是一个由“势—响应”坐标结构内生逼出的候选，而非外加猜测。
对该包络的一个可审计数值检验（通过 \(\nabla^2P(\theta)\) 的“单尺度/各向同性”程度扫描 \((u,v,w)=(e^{\theta_1},e^{\theta_2},e^{\theta_3})\)）见附录注 \ref{rem:sync-kernel-3d-sp6-isotropy-scan} 与表 \ref{tab:sync_kernel_weighted_sp6_isotropy_scan}。
\end{remark}

\paragraph{新结论 B（Zeckendorf--GUT 签名）：统一代数 \texorpdfstring{$\leftrightarrow$}{<->} 分辨率层组合}
对任意候选统一李代数 \(\mathfrak g\)，令 \(D:=\dim\mathfrak g\)（即规范玻色子总数）。
以 Zeckendorf 唯一分解写
\[
D=\sum_{k\in S}F_k,\qquad S\ \text{不含相邻指标},
\]
并定义其\textbf{分辨率签名}
\[
\Sigma(\mathfrak g):=\{\,m=k+2:\ k\in S\,\}.
\]
由于 \(F_k=|X^{\mathrm{bdry}}_{k+2}|\)，这等价于
\[
D=\sum_{m\in\Sigma(\mathfrak g)}|X_m^{\mathrm{bdry}}|,
\]
并且由 Zeckendorf 唯一性，\(\Sigma(\mathfrak g)\) 也是\emph{唯一的}（无冗余激活）分辨率层集合。
若把“可观测的 gauge 模态”理解为某些边界层 \(X_m^{\mathrm{bdry}}\) 在协议核/投影门下被激活的子塔，
则 \(\Sigma(\mathfrak g)\) 提供了一套可反驳的字典：只需在你的可审计协议中确定哪些 \(m\) 参与了可观测作用，即可把观测落入对应的统一签名。

\input{sections/generated/tab_zeckendorf_gut_signatures}

\paragraph{最小新预测（离散跳跃而非连续可调）}
一旦接受“规范玻色子 = 边界词塔模态数”的解释，SM 的签名为
\(\Sigma_{\mathrm{SM}}=\{4,6,8\}\)。
那么“超越 SM 的新增规范模态”不再是连续可调参数，而必须以 Fibonacci 边界层作离散跳跃：
最小的上跳新增层为 \(m=9\)（新增 \(F_7=13\) 模态）或 \(m=10\)（新增 \(F_8=21\) 模态）。
其中 \(m=10\) 的 \(21\) 模态允许一个与 \(SU(5)\) 重玻色子计数同型的 \(21-(8+1)=12\) 余量拆分，
为后续把“高层边界塔向低层三塔的投影/补偿”做成算子化闭合提供直接靶标。

\subsection{Fibonacci--Lie 共振梯：把规范群结构解释为分辨率窗口的离散共振}\label{subsec:fib-lie-resonance}

\paragraph{稳定类型维数的“非外加”离散锚点}
在黄金均值语法（禁止相邻 \(11\)）下，长度 \(m\) 的稳定类型集合满足
\(
|X_m|=F_{m+2}
\)。
因此当我们把“规范自由度的生成元数/玻色子数”视为某一层可被读出的离散维数接口时，
一个自然且可审计的候选是：在若干小窗口上，\(|X_m|\) 与紧李代数的维数发生精确相等，
从而把 \(\mathrm{SU}(2)\)、\(\mathrm{SU}(3)\) 等结构解释为分辨率增长过程中的\emph{共振点}，而非外加假设。

\paragraph{结论 C（硬结果）：Fibonacci--Lie 共振梯（可复现穷举）}
表 \ref{tab:fibonacci_lie_resonance_ladder} 列出在 \(m\le 500\) 的穷举扫描中出现的全部精确共振点：
\input{sections/generated/tab_fibonacci_lie_resonance_ladder}
并且更强地，若把等号直接写成二次丢番图方程，则在同一穷举范围内只出现极少数解：
\input{sections/generated/eq_fib_lie_resonance_solutions}
上述两段输出均由脚本 \texttt{scripts/exp\_fibonacci\_lie\_resonance\_ladder.py} 生成，
属于可审计、可复现的算术事实（不依赖任何解释层叙事）。

\begin{remark}[与边界词塔的四步平移]\label{rem:fib-lie-resonance-shift-by-4}
在第 \ref{subsec:bdry-tower-zeck-gut} 小节中，我们使用端点 \((1,1)\) 纤维的基数
\(
|X_m^{\mathrm{bdry}}|=F_{m-2}
\)
作为“边界模态数”。
注意到
\[
|X_m^{\mathrm{bdry}}|=F_{m-2}=F_{(m-4)+2}=|X_{m-4}|,
\]
因此本小节的共振梯（基于 \(|X_m|\)）在边界塔口径下表现为一个 \(m\mapsto m+4\) 的离散平移：
稳定类型共振点 \(m=2,4,6,8,\dots\) 对应边界层 \(m=6,8,10,12,\dots\)。
这解释了为何边界三层塔（\(m=4,6,8\)）会自然对齐 SM 的 \(1,3,8\) 三段玻色子计数接口：
它是同一 Fibonacci 维数骨架在不同纤维截面上的读出。
\end{remark}

\begin{remark}[最小额外生成元计数：\(21-12=9\)]\label{rem:fib-lie-extra-9}
若把 \(m=6\) 的稳定类型维数 \( |X_6|=21 \) 视为“最小统一层”的生成元数候选，
并把 SM 的可见玻色子数视为 \(12\)（第 \ref{subsec:bdry-tower-zeck-gut} 小节的边界塔接口），
则二者之间的差额为一个刚性的整数：
\[
21-12=9.
\]
这提供一个与“先选大群再分解”的先验群分解路线不同的计数指纹：额外通道必须以 \(9\) 个生成元的离散尺度出现，
其具体实现应通过本框架的扭曲谱/边界层离散标签/投影门机制来审计（而非作为连续可调参数引入）。
\end{remark}

\begin{remark}[fold-gauge 的 \texorpdfstring{$(\ZZ_2)^{21}$}{(Z2)^21} 二值电荷层（\texorpdfstring{$m=6$}{m=6}，dyadic 读出）]\label{rem:fold-gauge-z2-21}
除稳定类型维数 \(|X_6|=21\) 的“维数层”外，\(m=6\) 在纤维冗余层还给出一个独立、可审计的电荷格。
对二进制区间折叠 \(\mathrm{Fold}^{\mathrm{bin}}_6:\{0,\dots,63\}\to X_6\)（定义 \ref{def:fold-bin}），其纤维自同构群
\[
\mathcal{G}^{\mathrm{bin}}_6:=\mathrm{Aut}\!\bigl(\mathrm{Fold}^{\mathrm{bin}}_6\bigr)
\]
满足
\[
\bigl(\mathcal{G}^{\mathrm{bin}}_6\bigr)^{\mathrm{ab}}\cong (\ZZ_2)^{21}
\]
（推论 \ref{cor:fold-bin-gauge-m6}）。
这里每个 \(\ZZ_2\) 因子可由对应纤维内置换的符号表示读出，可解释为“最小非平凡窗口上可见的二值 gauge 电荷”。
该 \(21\) 维二值电荷层与 \S\ref{subsec:bdry-tower-zeck-gut} 的 \(B_3/C_3\) 种子字典在维数上吻合，并为进一步把
\(
21=\dim\wedge^2\RR^7=\dim\mathfrak{so}(7)
\)
这类几何解释升级为可审计约束提供一条额外证据链。
\end{remark}
\input{sections/08_group_unification_spectral_alignment.tex}\subsection{大统一方向的推进结论：几何化、跑动与可证伪锚点}\label{subsec:group-unification-gut-forward}

在 \S\ref{subsec:group-unification-spectral-alignment} 已将统一判据落实为谱隙耦合 \(g_{\varrho}(u)\) 的对齐与自对偶 completion 下统一点钉在 \(u=1\)（注 \ref{rem:selfdual-u1-six-branch-lock}），并给出数值锚点 \(W_{\mathrm{GUT}}\approx 5.32\)（命题 \ref{prop:real-input-40-unified-generator-alignment}）。本小节在上述结构之上给出六条可直接落地的推进结论及三条可检验预言，并在末尾补充七条接口级结论（以可检验形式陈述）；各条均回扣既有符号与接口，但不复述前文已写明的编号结论。

\paragraph{推进结论 1：统一点 \(u=1\) 的几何化为 \(\Omega\)--反演喉道固定点}
在极坐标全息动力学（holographic polar dynamics）的语境中，黑洞外各向同性坐标存在反演对称
\(
I:\ \rho\mapsto\tilde\rho=\rho_h^2/\rho
\)，
空间切片具有两渐近平坦端 \(\rho\to\infty\) 与 \(\rho\to 0\)，在 \(\rho=\rho_h\) 的 Einstein--Rosen 喉道处粘合。取无量纲尺度 \(u:=\rho/\rho_h\)，则 \(\Omega\)--反演为 \(u\mapsto 1/u\)，其唯一固定点为 \(u=1\)。

\begin{proposition}[统一点与喉道固定点的同一]\label{prop:u1-throat-identity}
自对偶 completion 下将统一点钉在 \(u=1\) 的谱学事实（注 \ref{rem:selfdual-u1-six-branch-lock}）可提升为几何同一：统一点 \(u=1\) 等同于 \(\Omega\)--反演喉道固定点（两张渐近平坦片的粘合处）。因此统一不是“某能标上耦合偶然相等”，而是协议在喉道处不得不执行的自对偶翻转（扫描/读出角色互换的几何化）；统一点固定而非参数可调，源于 \(u\mapsto 1/u\) 的强制对称。
\end{proposition}

\paragraph{推进结论 2：自对偶群提升为 Möbius/模群型对偶结构}
同一极坐标动力学中可引入 Cayley 型全息反演因子
\(
J(\rho)=(\rho-\rho_h)/(\rho+\rho_h)
\)，
满足反演下 \(J(\tilde\rho)=-J(\rho)\) 且 \(g_{tt}=-J^2\) 不变。在分辨率折叠的 \(\zeta\) 正规化中，以 \(z=r/\varphi\) 将主极点规范到 \(r=1\)，并将“允许禁转移”解释为单位盘内部的 pole-barrier 语言。

\begin{proposition}[三种自对偶/边界固定点的 Möbius 统一]\label{prop:mobius-duality-unified}
体系中已有的三种自对偶/边界固定点可纳入同一 Möbius 对偶结构：
\begin{itemize}
  \item \(\Omega\)--反演：\(u\mapsto 1/u\)（喉道固定点 \(u=1\)）；
  \item Cayley：\(J\mapsto -J\)（翻图不翻度量的对偶）；
  \item \(\zeta\) 的 pole-barrier：主极点钉到 \(r=1\) 的单位盘边界（可稳定延拓的边界）。
\end{itemize}
三者可在协议参数空间上由 Möbius 变换生成的对偶群组织；在最简生成元下与典型模群生成方式（反演与平移）同构。统一因而不仅是谱隙对齐，而是协议参数落在自对偶群的固定点/基本域边界上，为将 Ramanujan/模形式/素数轴与 \(\Omega\)--几何、谱隙跑动置于同一对偶图提供硬接口。
\end{proposition}

\paragraph{推进结论 3：谱隙跑动的 \(\beta\) 函数闭式（PF 期望）}
沿用定义 \ref{def:spec-coupling} 的谱隙耦合 \(g_{\varrho}(u)=\log(\lambda_1(u)/\eta_{\varrho}(u))\)。设 \(M(u)\) 为正矩阵族（或适用 Perron--Frobenius 理论的离散化），\(\rho(M(u))\) 为 Perron 根，左、右 PF 本征向量满足 \(M r=\rho r\)，\(\ell^\top M=\rho \ell^\top\)，\(\ell^\top r=1\)。则有
\[
\frac{\dd}{\dd u}\log\rho(M(u))
=\frac{\ell(u)^\top M'(u)\,r(u)}{\rho(M(u))}.
\]
定义协议 \(\beta\) 函数（以 \(\log u\) 为尺度）：
\[
\beta_\varrho(u)
:=\frac{\dd}{\dd\log u}\,g_\varrho(u)
=\frac{\dd}{\dd\log u}\log\lambda_1(u)
-\frac{u\,\ell(u)^\top M'(u)\,r(u)}{\rho(M(u))}.
\]

\begin{proposition}[\(\beta\) 的 PF 期望形式]\label{prop:beta-pf-expectation}
在本体系内，重整化群跑动无需微扰展开：\(\beta\) 函数等于“基准谱”与“通道 PF 压力”的两项可审计期望差，通道项为 \(M'(u)\) 在 PF 平衡态下的期望。这将“统一=谱隙对齐”的静态判据升级为统一点附近跑动由协议动力学锁定的动力学判据，且完全有限维、可复现。
\end{proposition}

\paragraph{推进结论 4：二阶 Jet 对齐作为稳定性选择器}
仅凭 \(g_{\chi,\varrho}(u)\) 在某点的值对齐不足以保证统一稳定：轻微改变观测门或核可能导致统一点漂移或“表面对齐、跑动分叉”。

\begin{proposition}[Jet 对齐的强统一判据]\label{prop:jet-alignment-unification}
真正可筛掉“伪统一”的判据为在统一点处的 Jet 对齐：
\[
g_{\chi,\varrho}(1)\ \text{对齐},\quad
g_{\chi,\varrho}'(1)\ \text{对齐},\quad
g_{\chi,\varrho}''(1)\ \text{对齐}.
\]
等价地，在统一点处所有激活通道的尺度势在各自 PF 平衡态下的一阶矩与二阶矩同时匹配。二阶导数可写为含协方差/二次响应项的表达式（压力函数的二阶响应）。因此可将“统一”从标量条件升级为可被数据拒绝的强条件：若某通道的曲率 \(g''(1)\) 不匹配，则判定其不属于同一统一类；同时匹配一阶与二阶响应可显著提高群选择器的锋利度。
\end{proposition}

\paragraph{推进结论 5：分辨率折叠与跑动的最小步长 \(\Delta_{\min}\)}
分辨率折叠给出平均纤维规模 \(d_m=2^m/F_{m+2}\) 及规范化极限 \(d_m(\varphi/2)^m\to\sqrt5/\varphi^2\)（\cite{Internal2025ResolutionFoldingPhiPiE}），故
\(
\log d_m=m\log(2/\varphi)+O(1).
\)
规范差大偏差接口（命题 \ref{prop:fold-gauge-anomaly-ldp}）给出端点稀有事件指数 \(I(0)=\log(2/\varphi)\)。

\begin{proposition}[\(\beta\) 差值的量子化步长]\label{prop:beta-quantization-step}
体系内存在由分辨率折叠与规范差端点代价共同锁定的不可约步长
\[
\Delta_{\min}:=\log(2/\varphi),
\]
它同时控制“折叠隐藏度的指数增长”与“规范差端点的大偏差代价”。因此所有通道的 \(\beta\) 函数在统一点附近的差值应呈现整数倍量子化：
\[
\beta_i(1)-\beta_j(1)\in \Delta_{\min}\,\ZZ,
\]
其中整数来源于通道在折叠尾部分布上占用的 Fibonacci 层数/权重。该形式在传统 RG 文献中非标准，但在本体系内自然：跑动由分辨率折叠导致的有效自由度变化给出，自由度变化的指数底数为 \(2/\varphi\)。
\end{proposition}

\paragraph{推进结论 6：\(W_{\mathrm{GUT}}\) 的闭式候选与可证伪性}
体系内必然出现的常数包括 \(\varphi\)（折叠/稳定语法/黄金分支）、\(\log(2/\varphi)=I(0)\)（规范差端点代价）以及几何常数 \(\pi\)。

\begin{conjecture}[\(W_{\mathrm{GUT}}\) 的闭式候选]\label{conj:Wgut-closed-form}
统一常数 \(W_{\mathrm{GUT}}\)（命题 \ref{prop:real-input-40-unified-generator-alignment}）可能由体系内已出现常数的小整数系数线性组合给出：
\[
W_{\mathrm{GUT}}\stackrel{?}{=}\pi+2\varphi-5\log(2/\varphi).
\]
右端数值约为 \(5.3179838535\ldots\)，与文中 \(5.32\) 在 \(10^{-3}\) 量级内一致；通过提高定义 \(W(v)\) 的脚本精度即可对该式作直接的数值证伪；若通过，则可在更高小数位上继续进行一致性检验。若成立，则得到非数秘的、可审计的“统一常数闭式”，系数可由通道数/层数/素数轴权重解释。
\end{conjecture}

\paragraph{可检验预言}
以下三条均可由现有脚本与谱数据直接检验。

\begin{enumerate}
  \item \textbf{喉道对偶预言（几何--谱）} 若将 \(u\) 解释为 \(\rho/\rho_h\)，则在统一类中应有近似对称
  \(
  g_{\chi,\varrho}(u)\approx g_{\chi,\varrho}(1/u)
  \)
  （因 \(\Omega\)--反演为 \(u\mapsto 1/u\) 且两端粘合）。检验方式：取对称点对 \((u,1/u)\)，检查谱隙耦合的镜像误差是否在证书界内。

  \item \textbf{\(\beta\) 差量子化预言} 在统一点附近，任意两通道的 \(\beta(1)\) 之差应落在 \(\log(2/\varphi)\) 的整数倍网格上。该条件比“耦合值接近”更苛刻，可排除仅在某点偶然对齐的模型。

  \item \textbf{Jet 对齐预言（强稳定性）} 若候选统一方案真实，则在不同观测门/核的微扰下，统一点 \(u=1\) 的二阶响应 \(g''(1)\) 应保持同类（误差随证书缩小）；否则说明统一由参数微调造成的偶然碰撞所致。
\end{enumerate}

\paragraph{补充结论与可检验预测（统一骨架的接口级表述）}\label{par:gut-seven-additional-results}
以下七条均以本文已建立的 \(G_\infty^{\mathrm{GUT}}=\widehat{\ZZ}\times G_{\mathrm{chan}}\)、提升算子 \(T_{\mathrm{GUT}}\) 与双对偶分解
\(
T_{\mathrm{GUT}}\simeq\bigoplus_{\chi}\bigoplus_{\varrho}M_{\chi,\varrho}
\)
为共同底座（注 \ref{rem:pom-grand-unified-lift}），并只给最终断言与可检验形式，不展开推导。
\begin{enumerate}
  \item \textbf{（补充结论 1：三通道耦合的单对象化——同一算术对象的三个 Satake 分量）}
  在“prime-skeleton + Euler/Hecke 闭包 + 可审计系数恢复链”的协议口径下，三条投影通道（注 \ref{rem:pom-grand-unified-lift} 的 \((\varphi,\pi,e)\)）在素数处的局部指纹不应被视为三条独立自由曲线；
  更强的统一形态应为：存在一个全局 Hecke 本征对象（等价地，一个全局 \(L\)-数据/自守表示候选）使得三通道的局部数据是同一局部因子的三个可计算分量（最自然候选为 \(GL(2)\) 本征对象的 \(\mathrm{Sym}^2\) 提升所诱导的三重局部参数）。
  \emph{可检验形式：}
  （i）对每个素数 \(p\)，三通道局部指纹必须满足同一族 Hecke/Euler 递推约束（不允许通道间相对漂移自由度）；
  （ii）全局一致性可用有限确定性模板完成：依 \S\ref{subsec:group-unification-audit-pointers} 的 R4（Hecke 刚性与 Sturm 界口径），比较有限素数集上的局部数据即可对“单对象化”假设作有限判定。

  \item \textbf{（补充结论 2：\(G_{\mathrm{chan}}\) 的最小闭合——\(U(3)\) 与内生 \(\ZZ_2\) 的半直积）}
  在连续包络口径下，三通道势—响应辛结构强迫最小连续对称为 \(\mathrm{Sp}(6,\RR)\)，其最大紧子群为 \(U(3)\)（注 \ref{rem:sp6-minimal-continuous-envelope}）。
  与此同时，window-6 的 dyadic uplift 在 boundary 扇区内生产生二层覆盖与 sheet parity
  \(\pi_{\mathrm{sheet}}\in\ZZ_2\)（\S\ref{subsec:bdry-tower-zeck-gut} 的结论 E）。
  因而 \(G_{\mathrm{chan}}\) 的最小闭合形态应收敛为
  \[
  \boxed{\;
  G_{\mathrm{chan}}\ \cong\ U(3)\rtimes \ZZ_2^{\mathrm{sheet}},
  \qquad \ZZ_2^{\mathrm{sheet}}=\langle \pi_{\mathrm{sheet}}\rangle,
  \;}
  \]
  其中离散 \(\ZZ_2\) 承载不可连续消去的“镜像/分支”超选择标签。
  \emph{可检验形式：}在统一类的可观测通道族上，应出现一个与 \(U(3)\) 混合自由度相对独立的二元分支；并且所有被允许的稳定更新/提升算子应与 \(\pi_{\mathrm{sheet}}\) 的分级相容（否则在审计证书上表现为跨分支的异常残差）。

  \item \textbf{（补充结论 3：统一耦合常数由 \(W_{\mathrm{GUT}}\) 数值锁定——\(\alpha_{\mathrm{U}}^{-1}=W_{\mathrm{GUT}}^2\)）}
  将真实输入 \(40\) 状态核上由方向耦合比 \(\mathcal{W}(v)\) 所读出的主导统一常数 \(W_{\mathrm{GUT}}\approx 5.32\)
  （命题 \ref{prop:real-input-40-unified-generator-alignment}）视作统一层的无量纲“刚度”锚点，则最简无参归一化给出统一耦合预测
  \[
  \boxed{\ \alpha_{\mathrm{U}}^{-1}:=W_{\mathrm{GUT}}^2\ \approx\ 28.30,\qquad
  g_{\mathrm{U}}:=\sqrt{\frac{4\pi}{\alpha_{\mathrm{U}}^{-1}}}=\frac{\sqrt{4\pi}}{W_{\mathrm{GUT}}}\ \approx\ 0.666\ }.
  \]
  \emph{可检验形式：}该预测不引入新数据，只依赖 \(W_{\mathrm{GUT}}\) 的复现精度；一旦 \(W_{\mathrm{GUT}}\) 的小数位被脚本证书锁定，上式即给出可直接检验的数值锚点。

  \item \textbf{（补充结论 4：统一阈值的离散量子化——按 \(\varphi^6\) 的六分支步进）}
  自对偶 completion 在统一点 \(u=1\) 处出现六分支锁定（注 \ref{rem:selfdual-u1-six-branch-lock}），提示统一层级的最小“尺度步进”应与 \(6\) 这一分支数绑定。
  在 \(\varphi\)-对数刻度
  \(
  r:=\log_{\varphi}(\mu/\mu_0)
  \)
  下，这给出一条离散阈值预测：统一阈值不连续可调，而按
  \[
  \boxed{\ \Delta r=6\qquad\Longleftrightarrow\qquad \frac{\mu_{m+1}}{\mu_m}=\varphi^6=9+4\sqrt5\approx 17.94427191\ }.
  \]
  \emph{可检验形式：}若把统一类的候选“阈值点/层级跃迁”落实为谱隙对齐或 adelic 平坦点的稳定极小族，则其在尺度轴上的出现位置应呈现 \(r\equiv 0\pmod 6\) 的离散层级结构（对核/门的微扰保持鲁棒，否则可被直接证伪）。

  \item \textbf{（补充结论 5：超选择禁止律——完成闭层中跨稳定指纹扇区的跃迁算符一阶缺失）}
  本文已给出一族在 \(p\)-进塔上稳定的审计不变量：\(\theta\)-导数的稳定首位残差 \(\alpha_{p,r}\in\mathbb{F}_p\)（推论 \ref{cor:witt-theta-stable-residue}）以及 \(p=2\) 的稳定奇偶指纹 \(\varepsilon_r\in\{0,1\}\)（推论 \ref{cor:witt-theta-binary-fingerprint}）；
  同时异常上同调类 \([\mathrm{Anom}(w)]\) 对水平复合严格加法（推论 \ref{cor:pom-anom-class-additive}）。
  因此在“完成闭层/稳定更新”口径下，允许算子代数应被这些稳定标签分解为直和的超选择扇区；任何必须改变 \(([\mathrm{Anom}],\{\alpha_{p,r}\},\{\varepsilon_r\})\) 的跃迁在一阶近似下不属于可实现算子。
  \emph{可检验形式：}给定任意候选“破缺算子/跃迁规则”，若其作用导致上述任一稳定指纹发生可审计改变，则可在有限层审计中直接否决其“闭层可实现性”。（在解释层中，若将某类物理量如“重子数改变”编码为超选择标签，则得到“主导衰变道一阶为零”的结构性预测。）

  \item \textbf{（补充结论 6：adelic anomaly 的强口径——乘积公式式的精确抵消与离散电荷量子化）}
  注 \ref{rem:pom-adelic-anom-functional} 已将局部异常强度 \(A_p(\theta)\) 汇总为整体泛函 \(\mathcal{A}_{\le P}(\theta)\) 并以稳定极小作为统一点判据。
  更强的统一口径是：异常不止“最小化”，而应在 adelic 粘合意义下形成全局平凡类（局部—整体精确抵消，类似乘积公式的闭合约束）。
  \emph{可检验形式：}随截断 \(P\to\infty\) 增大，局部异常签名的上同调类代表应在可审计规范下稳定，并且其全局类必须可被消去（否则统一点只能是截断依赖的伪极小）。
  该强口径同时带来一个硬后果：允许的离散电荷/超选择标签（例如 \(\pi_{\mathrm{sheet}}\) 及其与通道纤维的耦合方式）将被“可粘合性”条件一并量子化，从而消除一部分建模自由度。

  \item \textbf{（补充结论 7：有限资源律——可见统一表示必为截断显现，且 \(n_{\max}\sim N/\log N\)）}
  在任何以 Peter--Weyl/傅里叶分解读取 \((\chi,\varrho)\) 通道族的实现中，有限样本 \(N\) 的最坏频率误差由星差异度 \(D_N^*\) 确定性控制（\S\ref{sec:statistical-stability}；并见第 \ref{sec:experiments} 节的确定性 TV/KL 证书链）。
  因而存在一个不可回避的资源上界：可稳定恢复/判等的最大模式（表示）索引满足
  \[
  \boxed{\ n_{\max}\ \asymp\ \frac{1}{D_N^*}\ }.
  \]
  对 Kronecker/旋转扫描基准，当 \(D_N^*=O((\log N)/N)\) 时得到渐近资源律
  \[
  \boxed{\ n_{\max}\ =\Theta\!\left(\frac{N}{\log N}\right)\ }.
  \]
  \emph{可检验形式：}随着 \(N\) 增大，超过 \(n_{\max}\) 的高频模式应系统性地落入证书误差带（呈现“阶梯式显现”而非连续显现）；该阈值可由 \(D_N^*\) 的可计算上界直接预估并用于统一审计预算分配。
\end{enumerate}

\subsubsection{新增封闭结论：三通道同余提升的极小舞台与证书耦合}\label{subsubsec:gut-closed-conclusions}
在结论小节 \ref{subsubsec:bridgepack-e8-hecke-lock} 已将顶层统一母对象在“level-\(1\) 标量 + Hecke 刚性”口径下锁定为
\(
Z_{\mathrm{GUT}}(q)=E_4(q)=\Theta_{E_8}(q)
\)，
并给出“对称破缺的唯一出口”为提升到向量值/同余层、使 \((\varphi,\pi,e)\) 三投影成为三个自然分量（同小节之终结式结构结论 4）。
本段不复述上述硬输入的推导，而在其上记录六条可直接作为终点判别准则的封闭断言。

\paragraph{封闭结论 1：三通道分裂的极小同余舞台可取为 \texorpdfstring{$\Gamma_0(2)$}{Gamma0(2)} 的指数 \(3\) 覆盖，因而通道置换刚性为 \texorpdfstring{$S_3$}{S3}}\label{par:gut-closed-s3-index3}
若要求
(i)~顶层标量统一对象保持为 level-\(1\) 的 \(Z_{\mathrm{GUT}}\)（即不改变其 Hecke--素数壳骨架），并且
(ii)~分裂后通道数严格为 \(3\)（与 \((\varphi,\pi,e)\) 对齐），
则三通道结构可在同余几何上取为 \(\Gamma_0(2)\subset \mathrm{PSL}_2(\ZZ)\) 的指数-\(3\) 覆盖；对应的三通道表示可取为余类置换表示
\[
\rho_{\mathrm{chan}}\;\cong\;\mathrm{Ind}_{\Gamma_0(2)}^{\mathrm{PSL}_2(\ZZ)}\mathbf{1},
\]
从而三通道的内禀离散对称在同构意义下刚性等同于 \(S_3\)。
\emph{可检验形式：}任意候选向量值提升若声称“三通道分裂来自同余覆盖的极小性”，则其通道空间必须承载一个与 \(S_3\) 置换表示同构的离散对称；否则该解释不可成立。

\paragraph{补充封闭结论：\texorpdfstring{$V_4$}{V4} 通道核、\texorpdfstring{$N=8$}{N=8} 同余实现与模 \(8\) 符号律}\label{par:gut-closed-v4-mod8}
若进一步要求三通道在闭层对应三个彼此对易的 \(\ZZ_2\) 型超选择投影（从而通道可分解且不可任意混合），则通道荷群的极小有限闭包可取为 Klein 四元群
\[
V_4\cong \ZZ_2\times \ZZ_2.
\]
其特征标群 \(\widehat{V_4}\) 恰含三个非平凡一维特征标，从而可取
\[
(\varphi,\pi,e)\ \longleftrightarrow\ \widehat{V_4}\setminus\{\mathbf{1}\}
\]
作为通道识别。
该选择与封闭结论 1 的 \(S_3\) 刚性并不冲突：\(\mathrm{Aut}(V_4)\cong S_3\) 作用于三个非平凡特征标，提供一个把“通道置换”解释为“通道核自同构”的闭层对齐。
\par
进一步地，\(V_4\) 可内嵌为 \(U(3)\) 的对角符号翻转子群（例如由 \(\mathrm{diag}(-1,-1,1)\)、\(\mathrm{diag}(-1,1,-1)\)、\(\mathrm{diag}(1,-1,-1)\) 与单位元生成），因此与补充结论 2 的最小连续包络兼容。
\par
在算术实现上，取 \((\ZZ/8\ZZ)^\times\cong V_4\) 并用其三条非平凡实 Dirichlet 特征标
\(\chi_{-1},\chi_{2},\chi_{-2}=\chi_{-1}\chi_2\)
给出三通道的同余提升（极小层级 \(N=8\)）。
在素数壳审计口径下（良素数 \(p\neq 2\)），若以
\(
a_p=240(1+p^3)
\)
为 \(E_4\) 的既定素数系数，则三通道系数满足
\[
a_p^{(\chi)}=\chi(p)\,a_p,\qquad \chi\in\{\chi_{-1},\chi_2,\chi_{-2}\}.
\]
从而三通道的“素数壳”被压缩为仅依赖 \(p\bmod 8\) 的符号律；例如约定 \((\varphi,\pi,e)\equiv(\chi_2,\chi_{-1},\chi_{-2})\)，则对任意奇素数 \(p\) 有表 \ref{tab:gut-mod8-sign-law}。
\begin{table}[H]
\centering
\caption{模 \(8\) 符号律：三通道素数壳系数相对 \(E_4\) 的符号（奇素数）}
\label{tab:gut-mod8-sign-law}
\begin{tabular}{rrrrr}
\toprule
\(p\bmod 8\) & \(\chi_2(p)\) & \(\chi_{-1}(p)\) & \(\chi_{-2}(p)\) & \((a_p^\varphi,a_p^\pi,a_p^e)\) 相对 \(a_p\)\\
\midrule
\(1\) & \(+1\) & \(+1\) & \(+1\) & \((+,+,+)\)\\
\(3\) & \(-1\) & \(-1\) & \(+1\) & \((-,-,+)\)\\
\(5\) & \(-1\) & \(+1\) & \(-1\) & \((-,+,-)\)\\
\(7\) & \(+1\) & \(-1\) & \(-1\) & \((+,-,-)\)\\
\bottomrule
\end{tabular}
\end{table}
\emph{可证伪形式：}一旦从三通道读出中提取的素数壳系数在足够多的奇素数上不满足该模 \(8\) 符号律，则“\(G_{\mathrm{chan}}=V_4\) + level \(8\) 的最小 Dirichlet 扭曲实现”被直接否决。

\paragraph{封闭结论 2：三通道对象必分解为 \texorpdfstring{$\mathbf{1}\oplus\mathbf{2}$}{1+2}，因而“统一”成为严格守恒的全对称投影}\label{par:gut-closed-1plus2}
由 \(S_3\) 的三维置换表示结构可知必有
\[
\rho_{\mathrm{chan}}\;\cong\;\mathbf{1}\;\oplus\;\mathbf{2},
\]
其中 \(\mathbf{1}\) 为不变线、\(\mathbf{2}\) 为标准二维不可约分量。
因此对任意三分量提升
\(
Z^{\mathrm{vec}}_{\mathrm{GUT}}=(Z_\varphi,Z_\pi,Z_e)
\)，
其全对称投影（通道和/迹映射）必严格回落到标量子空间。
而在既定权重与 level-\(1\) 口径下该标量空间为一维，故存在常数 \(c\) 使
\[
Z_\varphi+Z_\pi+Z_e=c\cdot Z_{\mathrm{GUT}}.
\]
若采用常数项归一化（每通道 \(q^0\) 系数取 \(1\)），则 \(c=3\)。
\emph{可检验形式：}三通道分裂只能发生在二维标准分量内，因此三通道之间至多具有两条独立“差异自由度”；上述对称投影恒等式在任何候选实现中均为硬约束。

\paragraph{封闭结论 3：window-\(6\) 的端点三态给出三通道的规范基底，从而建立“有限窗语法—同余覆盖”之间的非任意同构}\label{par:gut-closed-bdry-basis}
在 window-\(6\) 锚点，端点纤维（boundary 扇区）的三条边界词为
\[
X^{\mathrm{bdry}}_6=\{\texttt{100001},\ \texttt{100101},\ \texttt{101001}\}
\qquad(\text{见 }\texttt{\detokenize{sections/generated/eq_fold6_b3c3_seed.tex}}).
\]
该三态可作为三通道指数集的规范基底，并与封闭结论 1 的指数-\(3\) 余类基底在同构意义下一一对应。
由此，通道标签 \((\varphi,\pi,e)\) 的“离散三性”可在闭层内同时由
（语法侧）window-\(6\) 的端点三态与
（算术几何侧）指数-\(3\) 的同余覆盖
两种独立构造一致地产生，从而把通道标签从外加语义降低为可审计的结构事实。
\emph{可检验形式：}任何候选向量值/同余提升若不能在其通道基底中（至少在同构意义下）重现该端点三态的规范性，则与有限窗闭层接口不相容。

\paragraph{封闭结论 4：通道混合的坏素数支撑律——非对角 Hecke 混合只能发生在同余层级的坏素数上}\label{par:gut-closed-bad-prime-support}
在封闭结论 1--2 的同余提升中，Hecke 作用在通道空间的\emph{非对角混合}只可能由同余层级 \(N\) 的坏素数（即 \(p\mid N\)）触发；
对所有良素数 \(p\nmid N\)，Hecke 作用在 \(\mathbf{1}\oplus\mathbf{2}\) 分解上可同时对角化，并保持不变线与标准分量的分离。
因此，若要求与本文已反复出现的三轴集合 \(\{2,3,5\}\) 相容（把其视为“通道混合的最小算术支撑集”），则同余层级必须满足 \(2,3,5\mid N\)，其最小候选为 \(N=30\)。
\emph{可证伪形式：}一旦在足够多的良素数处观测到不可消去的通道混合，则“层级 \(N\) 的同余提升”在该 \(N\) 下被直接否决；反之，混合若仅支撑在坏素数集合上，则其算术来源被压缩为有限坏素数组合。

\paragraph{封闭结论 5：在动机—Langlands 升级链上，三通道破缺对应一个次数 \(3\) 的置换型 Artin 因子}\label{par:gut-closed-artin-factor}
在封闭结论 1 的表述下，三通道结构
\(
\rho_{\mathrm{chan}}\cong \mathrm{Ind}_{\Gamma_0(2)}^{\mathrm{PSL}_2(\ZZ)}\mathbf{1}
\)
等价于一个次数 \(3\) 的置换型 Artin 表示/动机因子。
与顶层标量统一对象结合后，“破缺”可被写成
\[
(\text{顶层 }Z_{\mathrm{GUT}}\text{ 的 Eisenstein/}\theta\text{ 型骨架})\;\otimes\;(\text{次数 \(3\) 的 Artin 因子})
\]
的组合，从而把破缺机制外推为可分类的有限对象（选择同余层级、有限表示与坏素数处的局部矩阵数据），并与本文强调的“零参数 Hecke 骨架”审计目标兼容。

\paragraph{封闭结论 6：统一与 RH/GRH 的证书化在同一闭层审计框架下发生耦合（统一证书 \texorpdfstring{$\Leftrightarrow$}{<->} \texorpdfstring{$\zeta$}{zeta}-因子证书的可兼容性）}\label{par:gut-closed-certificate-coupling}
当顶层统一对象属于 Eisenstein/\(\theta\)-型并通过 Hecke 数据确定其 \(L\)-因子骨架时，统一闭包所需的“全尺度一致性证书族”（保证跨尺度读出稳定、排除不可控漂移）在数论层面必然要求相应 \(\zeta\)-类因子满足可兼容的零点/谱性质证书。
因此在该路线中：
“统一闭包是否可完成”与“相关 \(\zeta\)-因子（及其平移/乘积）的 RH/GRH 型证书是否可被同一审计机制承载”
不再是两条独立问题，而成为同一证书体系的可组合性兼容性约束。

\paragraph{三条终局判别标准（无需过程性中间结论）}\label{par:gut-closed-three-criteria}
\begin{enumerate}
  \item \textbf{同余极小性判据：}任何三通道候选必须在同余/向量值层面可等价实现为指数-\(3\) 的覆盖诱导结构（封闭结论 1），否则无法解释“通道数 \(=3\)”的结构极小性。
  \item \textbf{守恒投影判据：}必须存在一个通道全对称投影使其严格回落到 level-\(1\) 的标量 \(Z_{\mathrm{GUT}}\)（封闭结论 2），即 \(Z_\varphi+Z_\pi+Z_e\propto Z_{\mathrm{GUT}}\)。
  \item \textbf{坏素数混合判据：}任何非对角通道混合只允许出现在坏素数集合 \(p\mid N\) 上；若需要与三轴集合 \(\{2,3,5\}\) 相容，则坏素数集合至少包含 \(\{2,3,5\}\)，从而 \(30\mid N\)（封闭结论 4）。
\end{enumerate}

\paragraph{下一步信息增益唯一点（从断言到确定对象）}
在封闭结论 1--4 的约束下，进一步推进到“可提交的确定对象”的唯一信息增益工作是：
构造并固定一个具体的三分量 \(Z^{\mathrm{vec}}_{\mathrm{GUT}}\)（含其 Hecke 算子在坏素数处的矩阵数据与在良素数处的本征分解），
并用封闭结论 2 的守恒投影将其与顶层 \(Z_{\mathrm{GUT}}\) 做迹一致性闭包；
该步骤一旦完成，即得到一个从顶层统一到通道分裂的零任意度蓝图。
\subsection{审计落脚点：统一命题的可复现接口}\label{subsec:group-unification-audit-pointers}

本节不引入新的证明前提，而是将本章中关于统一的关键结论逐条落实到仓库内的可复现输出与可证伪检查上，
从而读者无需诉诸解释层即可定位相应脚本、生成片段与审计证书。

\paragraph{可复核算术事实：Fibonacci--Lie 共振梯（可复现穷举）}
稳定类型计数律 \(|X_m|=F_{m+2}\)（本章口径）在小窗口上出现对紧李代数维数的精确共振（小节 \ref{subsec:fib-lie-resonance}）。
其关键部分是将等式写成丢番图方程并作有限窗口穷举，从而得到可复现的解集（与解释层叙述无关）：
\begin{itemize}
  \item \textbf{生成脚本：}\texttt{\detokenize{scripts/exp_fibonacci_lie_resonance_ladder.py}}。
  \item \textbf{LaTeX 生成物：}\texttt{\detokenize{sections/generated/tab_fibonacci_lie_resonance_ladder.tex}}，
  \texttt{\detokenize{sections/generated/eq_fib_lie_resonance_solutions.tex}}。
  \item \textbf{审计 JSON：}\texttt{\detokenize{artifacts/export/fibonacci_lie_resonance_ladder.json}}。
\end{itemize}
任何对共振点/解集的修改，都必须能通过上述脚本在声明的扫描窗口上复现。

\paragraph{可复核接口：\(m=6\) 的 \(18\oplus 3\) 分裂与 \(B_3/C_3\) 统一种子}
在边界词塔口径下，\(m=6\) 的端点纤维给出三条边界词（见 \texttt{\detokenize{sections/generated/eq_boundary_tower_sm12.tex}}），并由 \(X_6=X_6^{\mathrm{cyc}}\sqcup X_6^{\mathrm{bdry}}\) 的 \(18\oplus 3\) 分裂将 \(21\) 维统一种子限定为 \(B_3/C_3\)（小节 \ref{subsec:bdry-tower-zeck-gut}）。
为避免人工配对，本文将该种子进一步落实为根字典与 Weyl 群轨道大小的核验输出：
\begin{itemize}
  \item \textbf{生成脚本：}\texttt{\detokenize{scripts/exp_fold6_b3c3_seed.py}}。
  \item \textbf{LaTeX 生成物：}\texttt{\detokenize{sections/generated/eq_fold6_b3c3_seed.tex}}，
  \texttt{\detokenize{sections/generated/tab_fold6_b3c3_root_dictionary_B3.tex}}，
  \texttt{\detokenize{sections/generated/tab_fold6_b3c3_root_dictionary_C3.tex}}。
  \item \textbf{审计 JSON：}\texttt{\detokenize{artifacts/export/fold6_b3c3_seed.json}}（含 \(18\oplus 3\)、\(6+12\) 指纹、两套字典与 \(|W|=48\) 的轨道核验直方图）。
\end{itemize}

\paragraph{并列可复核接口：\(m=6\) 的 \texorpdfstring{$C_6$}{C6} 轨道分解与 \texorpdfstring{$21=15\oplus 3\oplus 3$}{21=15+3+3} Levi 骨架候选}
在同一 window-6 锚点上，环旋转 \(C_6\) 在 cyclic 扇区 \(X_6^{\mathrm{cyc}}\) 上封闭并诱导一个刚性的轨道多重集合 \(\{1,2,3,6,6\}\)（见小节 \ref{subsec:bdry-tower-zeck-gut} 中关于 \texorpdfstring{$C_6$}{C6} 轨道刚性的段落）。
该刚性轨道分解给出一个规范的 \(15\oplus 3\) 细分，并与 boundary 扇区的 \(3\) 一起形成 \(21=15\oplus 3\oplus 3\) 的维数骨架，
从而把一个 Pati--Salam 型的最小半单 Levi 包络作为并列候选落到可审计输出上。
此外，dyadic 折叠 \(\mathrm{Fold}^{\mathrm{bin}}_6\) 在 boundary 扇区上强迫 \(\delta=34\) 的二层覆盖并诱导一个可复核的 \(\ZZ_2\) sheet parity（见小节 \ref{subsec:bdry-tower-zeck-gut} 中 boundary 扇区双层覆盖的构造）。
\begin{itemize}
  \item \textbf{生成脚本：}\texttt{\detokenize{scripts/exp_window6_c6_orbit_patisalam_seed.py}}。
  \item \textbf{LaTeX 生成物：}\texttt{\detokenize{sections/generated/eq_window6_c6_orbit_decomposition.tex}}，
  \texttt{\detokenize{sections/generated/eq_window6_c6_character_decomposition.tex}}，
  \texttt{\detokenize{sections/generated/tab_window6_c6_orbit_decomposition.tex}}，
  \texttt{\detokenize{sections/generated/tab_window6_patisalam_9633_partition.tex}}，
  \texttt{\detokenize{sections/generated/eq_fold6_bin_uplift_delta_set.tex}}，
  \texttt{\detokenize{sections/generated/tab_fold6_bin_uplift_choice_collapse.tex}}，
  \texttt{\detokenize{sections/generated/tab_fold6_boundary_sheet_pairs.tex}}，
  \texttt{\detokenize{sections/generated/tab_foldbin_boundary_lift_m6_m8.tex}}。
  \item \textbf{审计 JSON：}\texttt{\detokenize{artifacts/export/window6_c6_orbit_patisalam_seed.json}}（含 \(C_6\) 轨道、\(9+6+3+3\) 四块分解、\(\Delta_6=\{0,21,34,55\}\)、uplift-choice collapse 的 \(\Xi(w)\in\{2,3,4\}\) 分区、\(w_6=1\Rightarrow\{0,34\}\)、三条边界词的显式 preimage 对，以及 \(m=6,7,8\) 的 boundary-lift 三层/二层稳定性扫描输出）。
\end{itemize}

\paragraph{可复核接口：\texorpdfstring{$X_6$}{X6} 的 \(1+8+12\) 字典与 \(12=8+4\) 两档重谱}
“\(1+8+12\)” 与 “\(12=8+4\)” 的两条离散分块结论无需另写脚本：
它们可直接从表 \ref{tab:fold6_bin_uplift_choice_collapse} 的三行词表读出。
第一行给出 \(S_{\mathrm{low}}\) 的 \(9\) 词（其中 \(\texttt{000000}\) 为唯一零词、其余 \(8\) 词构成 \(1+8\) 的轻块），
第二行与第三行合并为 \(12\) 的重块，并按 \(w_6=1\)（\(\Xi=2\)，\(8\) 词）与 \(w_6=0,\ V_6(w)>8\)（\(\Xi=3\)，\(4\) 词）自动分成两档。
由此可直接计算熵权重 \(h(w)=\log\Xi(w)\) 并作为谱隙对齐的分块权重（\S\ref{subsec:group-unification-spectral-alignment}）。
对应的可复核输出与脚本沿用上文关于 \(C_6\) 轨道与 dyadic uplift 证书的生成链（表 \ref{tab:fold6_bin_uplift_choice_collapse} 与
\texttt{\detokenize{artifacts/export/window6_c6_orbit_patisalam_seed.json}}）。

\paragraph{可复核接口：bin-fold 退化谱的双尺度选择律与 \texorpdfstring{$\dim=55$}{dim=55} 统一锚点}
bin-fold 退化谱的双尺度选择律将二进制区间折叠 \(\mathrm{Fold}^{\mathrm{bin}}_m\) 的退化度谱
\(
d^{\mathrm{bin}}_m(w)=|\mathrm{Fold}^{\mathrm{bin}}_m{}^{-1}(w)|
\)
作为独立于 Zeckendorf--GUT 签名与 NAP 选择器的第三条审计通道：
在 \(m=6\) 上由低退化扇区读出 \(12\) 的计数接口，并以 ghost 容量 \(H^{\mathrm{bin}}_6=2^6-|X_6|=43\) 补全得到 \(55\) 的统一维数锚点；
在更高分辨率 \(m=10\) 上，最小退化扇区计数再次精确等于 \(55\)，从而形成两路锁定。
其生成链已脚本化并纳入一键流水线：
\begin{itemize}
  \item \textbf{生成脚本：}\texttt{\detokenize{scripts/exp_fold_bin_degeneracy_spectrum_m6_m12.py}}。
  \item \textbf{LaTeX 生成物：}\texttt{\detokenize{sections/generated/tab_fold_bin_degeneracy_spectrum_m6_m12.tex}}（表 \ref{tab:fold_bin_degeneracy_spectrum_m6_m12}）。
  \item \textbf{审计 JSON：}\texttt{\detokenize{artifacts/export/fold_bin_degeneracy_spectrum_m6_m12.json}}（含各 \(m\) 的直方图、\(d_{\min}\)、最小扇区大小与墙钟时间）。
\end{itemize}

\paragraph{交叉审计接口：\texorpdfstring{$\mathrm{Fold}^{\mathrm{bin}}$}{Foldbin}–\texorpdfstring{$\mathrm{Fold}$}{Fold} 的跨尺度补余恒等式}
接口化结论 \ref{par:gut-foldbin-fold-complement-identity} 给出一条把“低分辨率最刚性扇区补余容量”与“高分辨率最大纤维极值模尺度”直接锁定的交叉审计恒等式；
其核对式为独立的派生输出（不重新枚举谱数据），并已纳入一键流水线：
\begin{itemize}
  \item \textbf{生成脚本：}\texttt{\detokenize{scripts/exp_foldbin_fold_complement_identity_m6_m12.py}}。
  \item \textbf{LaTeX 生成物：}\texttt{\detokenize{sections/generated/eq_foldbin_fold_complement_identity_m6_m12.tex}}。
  \item \textbf{审计 JSON：}\texttt{\detokenize{artifacts/export/foldbin_fold_complement_identity_m6_m12.json}}（含逐 \(m\) 核对项与通过/失败标记）。
\end{itemize}

\paragraph{可复核接口：四步平移 uplift 同构与 \texorpdfstring{$m=10$}{m=10} 的 \(6+6+6+3\) 根型分块}
四步平移 uplift 同构把 \(X_{10}^{\mathrm{bdry}}\) 与 \(X_6\) 规范同构，并以 \(B_3\) 根字典细化为
\(6\) 个短根、\(6\) 个 \(A_2(+)\)、\(6\) 个 \(A_2(-)\) 与 \(3\) 个 Cartan 方向。
对应的 \(21\) 词 uplift 字典与分块清单由脚本自动生成：
\begin{itemize}
  \item \textbf{生成脚本：}\texttt{\detokenize{scripts/exp_boundary_uplift_m10_b3c3_dictionary.py}}。
  \item \textbf{LaTeX 生成物：}\texttt{\detokenize{sections/generated/eq_boundary_uplift_m10_isomorphism.tex}}，
  \texttt{\detokenize{sections/generated/tab_boundary_uplift_m10_dictionary.tex}}。
  \item \textbf{审计 JSON：}\texttt{\detokenize{artifacts/export/boundary_uplift_m10_b3c3_dictionary.json}}（含 \(\iota_n\) 的有限窗双射核验、\(m=10\) 的 21 词字典与 \(A_2(\pm)\) 分块清单）。
\end{itemize}

\paragraph{可复核接口：NAP 选择器与 \texorpdfstring{$m=11$}{m=11} 的 \texorpdfstring{$\ZZ_{34}$}{Z34} Fourier 分解证书}
小节 \ref{subsec:bdry-tower-zeck-gut} 在 Zeckendorf--GUT 签名字典之上引入了一个新的离散选择律：
要求非阿贝尔共振层 \(\{6,8\}\) 必须作为独立层保留（NAP），并在所有简单李代数中最小化 \(D=\dim\mathfrak g\)。
该选择器把 \(SO(10)\) 作为最小解离散地挑出，并给出“\(4\to 11\)”的单次签名替换与 \(34-1=33\) 的刚性计数接口。
与此同时，\(m=11\) 的 \(34\) 个边界模态在规范编号下携带 \(\ZZ_{34}\) 平移作用，从而强制一个
\(34=1\oplus 33\) 且 \(33=1\oplus 16\times 2\) 的实表示分解模板（为后续投影门/核谱的谱并群提供可审计靶标）。
\begin{itemize}
  \item \textbf{NAP 选择器脚本：}\texttt{\detokenize{scripts/exp_nap_selector_so10.py}}。
  \item \textbf{NAP LaTeX 生成物：}\texttt{\detokenize{sections/generated/tab_nap_selector_so10.tex}}，
  \texttt{\detokenize{sections/generated/eq_nap_selector_so10.tex}}。
  \item \textbf{NAP 审计 JSON：}\texttt{\detokenize{artifacts/export/nap_selector_so10.json}}（含扫描参数、候选表与最小解证书）。
  \item \textbf{\(\ZZ_{34}\) Fourier 分解脚本：}\texttt{\detokenize{scripts/exp_m11_boundary_z34_fourier_decomposition.py}}。
  \item \textbf{\(\ZZ_{34}\) LaTeX 生成物：}\texttt{\detokenize{sections/generated/eq_m11_boundary_z34_decomposition.tex}}。
  \item \textbf{\(\ZZ_{34}\) 审计 JSON：}\texttt{\detokenize{artifacts/export/m11_boundary_z34_fourier_decomposition.json}}（含规范编号 \(j\mapsto (v_7,u_{11})\) 与正交/旋转校验残差）。
\end{itemize}

\paragraph{可复核检验：三通道势—响应辛相空间的 \texorpdfstring{$\mathrm{Sp}(6)$}{Sp(6)} 各向同性点扫描}
注 \ref{rem:sp6-minimal-continuous-envelope} 给出一个把 \(B_3/C_3\) 的离散统一种子“提升为连续包络”的最小口径：在压力 \(P(\theta)\) 的 Legendre 对偶坐标下，
\((\theta;J=\nabla P)\) 诱导标准辛形式，保持三通道共轭结构的最小连续群为 \(\mathrm{Sp}(6)\)。
对应的一个硬核可核对任务是：在某个“统一点”上，二次涨落尺度应趋于单尺度（各向同性），其完全可审计的分数为
\(
\mathcal{I}(\theta)=\log(\lambda_{\max}/\lambda_{\min})
\)
（见附录注 \ref{rem:sync-kernel-3d-sp6-isotropy-scan}）。
该扫描已被脚本化并纳入一键流水线：
\begin{itemize}
  \item \textbf{生成脚本：}\texttt{\detokenize{scripts/exp_sync_kernel_weighted_sp6_isotropy_scan.py}}。
  \item \textbf{LaTeX 生成物：}\texttt{\detokenize{sections/generated/tab_sync_kernel_weighted_sp6_isotropy_scan.tex}}（表 \ref{tab:sync_kernel_weighted_sp6_isotropy_scan}）。
  \item \textbf{审计 JSON：}\texttt{\detokenize{artifacts/export/sync_kernel_weighted_sp6_isotropy_scan.json}}（含扫描网格、\(\Sigma=\nabla^2P\)、特征值与最优候选点列表）。
  \item \textbf{根型分裂比率脚本：}\texttt{\detokenize{scripts/exp_sp6_isotropy_root_split_ratio.py}}（读取上条 JSON 并计算 \(R_{ls},R_{pm}\)）。
  \item \textbf{根型分裂 LaTeX 生成物：}\texttt{\detokenize{sections/generated/tab_sp6_isotropy_root_split_ratio.tex}}。
  \item \textbf{根型分裂审计 JSON：}\texttt{\detokenize{artifacts/export/sp6_isotropy_root_split_ratio.json}}。
\end{itemize}

\paragraph{接口延拓：将外部时间缩放、折叠缺陷与耦合漂移归约到同一主标量变量}
本章在注 \ref{rem:lapse-modulated-tau-mix} 中已将外部时间重标定写成一个最小的外部时间缩放代理
\(
N:=\kappa_0/\kappa
\)
（其中 \(\kappa\) 是背景依赖的路由/逆搜索开销，单位为“每内部步的外部代价”）。
另一方面，折叠侧已经定义了纤维缺陷标量
\begin{equation}\label{eq:gut_kappa_m_defect}
\kappa_m(\pi):=\mathbb{E}_{\pi}[\log d_m(X)]
\qquad(\text{定义 \ref{def:pom-kappa}；其中 }d_m(x)=|\Fold_m^{-1}(x)|).
\end{equation}
因此可以引入一个统一的主标量变量（纯定义，不引入额外指称）：
\[
u:=\log\frac{\kappa}{\kappa_0}=-\log N.
\]
\emph{接口级}假设是：在带空间索引的读出（例如 Hilbert 地址化窗口）下，局部开销场 \(u(x)\) 与某个局部折叠统计的缺陷量同构，
例如以局部分布 \(\pi_{m,x}\in\Delta(X_m)\) 读出
\(
\kappa_m(\pi_{m,x})
\)
并设
\[
u(x)\ \approx\ c\,\kappa_m(\pi_{m,x})
\qquad(c>0\ \text{为标定常数}).
\]
一旦将任意有效耦合或有效常数的漂移写成主标量的函数（线性化形式）
\[
g_i^{-2}(x)=a_i+b_i\,u(x)+\cdots,
\]
就得到一个\textbf{可证伪的相关性}：
\[
\Delta g_i^{-2}\ \approx\ b_i\,\Delta u\ =\ -b_i\,\Delta\log N.
\]
该相关性本身不依赖连续极限；它只要求在同一数据/同一协议族上同时输出 \(N\) 与 \(g_i\) 的代理量并做联合检验。
\par
更完整的 overhead\(\to\)外部时间缩放\(\to\)线性化闭合，以及从地址化窗口与折叠退化统计重建空间代价标量变量的可执行协议，属于接口层的进一步构造；
本稿在此仅记录其在本章变量上的最小同构关系与可证伪相关性形式。

\paragraph{有限检验模板：Hecke 刚性与 Sturm 界的有限素数骨架判定}
\begin{theorem}[审计化 Langlands 型识别的有限可判定性：Sturm 有界性 \(\times\) finite-layer closure \(\times\) counterterm strictification]\label{thm:audit-finite-decidability-sturm-finite-layer}
设某一候选“统一对象”的读出被组织为一组可比的系数/局部因子数据（例如模形式的 \(q\)-展开系数、Hecke 本征值 \(a_p\)，或等价的素数局部因子/特征多项式），并且其实现链允许同时输出 Dirichlet 扭曲通道与有限部漂移预算。
则在本文的可审计口径下，Langlands 型“识别/匹配”（是否属于给定的 Hecke 本征包/等价类）可归约为\textbf{有限}核验问题，其有限性来源可明确分解为三条：
\begin{enumerate}
  \item \textbf{（模形式端：Sturm 有界性）}在给定权 \(k\) 与水平 \(N\) 后，Sturm 界 \(B_{k,N}\) 给出有限截断判定：只需核验至 \(B_{k,N}\) 的有限个 Fourier/Hecke 数据即可判定等同性（从而把“无限条件”替换为有限素数骨架的检查）。
  \item \textbf{（Dirichlet 端：有限层闭包）}任何有限投影词/有限实现对 Dirichlet 轴束的依赖必在某个有限层 \(m_\ast\) 上闭合；因而“全角色扫描/最坏通道”可归约为对有限集合 \(\widehat{G_{m_\ast}}\) 的穷举核验（命题 \ref{prop:pom-finite-modulus-stability}，亦见推论 \ref{cor:pom-local-finite-dirichlet-action}）。
  \item \textbf{（接口严格化端：counterterm 消去有限部漂移）}中心扩张口径下的交换失败完全落在有限部/常数层，并可由纯常数 counterterm 门严格消去，从而把残差交换提升为严格相容性条件（命题 \ref{prop:pom-counterterm-anom-cancel}，亦见定理 \ref{thm:pom-strictification-splitting}）。
\end{enumerate}
因此，在既定的误差预算结构（有限样本/截断/strictification）下，上述识别问题不再是开放式的无穷验证，而是一个由 \((B_{k,N},m_\ast)\) 与有限部 counterterm 预算共同控制规模的有限判定问题。
\end{theorem}

\begin{remark}[总结性结论：lapse--Dirichlet--Adams--counterterm--Sturm 的闭合主线]\label{rem:audit-summary-lapse-dirichlet-adams-sturm}
本稿把“可审计计算”系统性地压缩为有限证书链：lapse（测度变换/外部时间缩放）可内化为加权周期读出并给出显式有限样本上界（定理 \ref{thm:op-algebra-lapse-weighted-wish-koksma}）；Dirichlet 扭曲通道在有限层闭包下局部有限并可穷举核验（命题 \ref{prop:pom-finite-modulus-stability} 与推论 \ref{cor:pom-local-finite-dirichlet-action}）；Adams--Dirichlet 的算术伸缩门把 \((\chi,u)\) 的幂半群作用统一为重整化单子（定义 \ref{def:witt-adams-dirichlet-renorm-gate} 与命题 \ref{prop:witt-renorm-monoid}）；异常通过中心扩张分裂与 counterterm strictification 被严格化为常数层可控漂移（定理 \ref{thm:pom-strictification-splitting}）。
在此基础上，Sturm 有界性把 Hecke 刚性落实为有限素数骨架判定（定理 \ref{thm:audit-finite-decidability-sturm-finite-layer}），从而形成一条无需手工编号、可逐条引用的闭合结构主线。
现有脚本 \texttt{\detokenize{scripts/exp_hecke_e4_prime_audit.py}} 即提供该模板的最小可复核实例：在零参数基准情形下，以有限素数点的 Hecke 系数表/JSON 输出作为快速否定护栏。
\end{remark}

\paragraph{零参数快速否定判据：权 \texorpdfstring{$4$}{4} Eisenstein/\texorpdfstring{$E_8$}{E8} \texorpdfstring{$\theta$}{theta}-模的素数骨架审计}
在上文“有限素数骨架判定”模板之上，一个特别刚性的基准情形是：若顶层统一对象被声明为
\emph{level-1（\(X(1)\)）标量权 \(4\) Hecke-闭合}的模形式，则其空间为一维并由 Eisenstein 级数 \(E_4\) 张成；
因此素数系数必须满足
\(
a_p=240(1+p^3)
\)
（\(p\) 为素数）。
\par
我们将最小否定判据素数点 \(p\in\{2,3,5\}\) 的数值表与 JSON 封装为可复现输出（作为“零拟合必要条件”的快速否定护栏）：
\begin{itemize}
  \item \textbf{生成脚本：}\texttt{\detokenize{scripts/exp_hecke_e4_prime_audit.py}}。
  \item \textbf{LaTeX 生成物：}\texttt{\detokenize{sections/generated/tab_hecke_e4_prime_audit_p235.tex}}。
  \item \textbf{审计 JSON：}\texttt{\detokenize{artifacts/export/hecke_e4_prime_audit_p235.json}}。
\end{itemize}
若扫描/读出得到的候选素数骨架 \(a_p\) 在这些素数点显著偏离，则说明顶层统一不可能仍停留在 “\(X(1)\) 标量 + Hecke 刚性” 这一层级，
必须进入向量值/同余层提升分支（即不再可能保持 \(X(1)\) 标量一维情形）。
